\chapter*{Exercices}
\addcontentsline{toc}{chapter}{Exercices}
\markboth{Exercices}{}

\subsection*{Exercice 1}

La répartition de l'énergie rayonnée par une étoile en fonction de la
longueur d'onde est donnée par la loi de Planck
\[
  I(\lambda) = (8\pi hc)\Big/\left(\lambda^5\left[\exp\Bigl(\frac{hc}{\lambda kT}\Bigr) - 1\right]\right)
\]
où $h$, $k$ et $c$ sont respectivement les constantes de Planck, de Boltzmann
et la célérité de la lumière. On pose $x = hc/\lambda kT$. D'après les
résultats expérimentaux, $I$ est maximale pour $x_m$.
\begin{enumerate}[label=\arabic*-]
  \item Établir l'équation vérifiée par $x_m$.
  \item On veut trouver $x_m$ par la méthode de Newton-Raphson. Établir
        l'algorithme permettant de calculer la valeur numérique de $x_m$.
\end{enumerate}

\subsection*{Exercice 2}

L'équation de Van der Waals pour une mole de gaz de masse $M$ est
\[
  \Bigl(P + \frac{a}{V^2}\Bigr)(V - b) = RMT
\]
où $a$, $b$ et $R$ sont des constantes. $P$, $T$ et $V$ sont respectivement la
pression, la température et le volume.

Établir un algorithme qui permet de calculer le volume $V$ chaque fois que
l'on connaît la température et la pression.

\subsection*{Exercice 3}

En 1976, Desantis a établi que le facteur de compressibilité $b$ des gaz réels
vérifie la relation
\[
  z = \bigl(1 + y + y^2 - y^3\bigr)\big/\bigl(1 - y^3\bigr)
\]
où $y = b/4v$ ; $v$ est le volume molaire et $z$ est une constante. On veut
trouver $b$ par la méthode de Newton-Raphson.
\begin{enumerate}[label=\arabic*-]
  \item Expliquer le principe de la méthode.
  \item Établir alors un algorithme permettant de résoudre le problème.
  \item Traduire cet algorithme en Fortran.
\end{enumerate}

\subsection*{Exercice 4}

Une grandeur $y$ vérifie l'équation suivante :
\[
  \frac{\dd y}{\dd x} = -ay - by^4 + c
  \qquad\text{avec}\quad y(0) = y_0 .
\]
$a$, $b$ et $c$ sont des constantes.

On se propose de résoudre cette équation numériquement par la méthode de
prédicteur-correcteur (Adams-Bashforth du second ordre associé à la méthode
de Runge-Kutta d'ordre 2).
\begin{enumerate}[label=\arabic*)-]
  \item Établir le schéma itératif correspondant.
  \item Établir l'algorithme de résolution de l'équation.
  \item Traduire cet algorithme en Pascal et en Fortran.
\end{enumerate}

\subsection*{Exercice 5}

Dans le domaine $x \in\, ]0, 1[$ une grandeur $y$ vérifie l'équation
différentielle
\[
  y'' + \frac{a}{x}y' + be^{cy} = 0
  \qquad\text{avec}\quad y'(0) = 0 \quad\text{et}\quad y(1) = 1 .
\]
$a$, $b$ et $c$ sont des constantes. On veut résoudre cette équation par la
méthode de tir (associée à la méthode de Runge-Kutta d'ordre 4).
\begin{enumerate}[label=\arabic*)-]
  \item En utilisant la méthode de Newton (pour les équations algébriques non
        linéaires), donner le schéma itératif qui permet de calculer les
        valeurs successives de la condition de tir en $x = 0$.
  \item Établir la méthode itérative de Runge-Kutta d'ordre 4.
  \item Établir l'algorithme complet (incluant le schéma des conditions
        initiales) permettant de résoudre le problème.
\end{enumerate}

\subsection*{Exercice 6}

Dans l'espace $KA$ entre une cathode et une anode $A$ parallèle à $K$, le
potentiel à la distance $x$ de $K$ satisfait l'équation différentielle
\begin{equation}
  V'' = aV^{-1/2}
  \tag{1}
\end{equation}
où $a$ est une constante. On se propose de trouver numériquement les valeurs
de $V$.
\begin{enumerate}[label=\arabic*-]
  \item Dans cette première question, on suppose que le potentiel et le champ
        électrique sont nuls sur la cathode (problème à valeurs initiales). On
        résout alors (1) par la méthode de Runge-Kutta d'ordre 2.
        \begin{enumerate}[label=\alph*-]
          \item Sachant que RK2 dérive de la formule intégrale des trapèzes,
                donner l'expression générale des itérations RK2 pour une
                équation différentielle du premier ordre.
          \item L'équation (1) peut se décomposer en deux équations
                différentielles du premier ordre. Établir l'algorithme de
                résolution numérique de (1) par RK2.
          \item Traduire cet algorithme en Pascal et en Fortran.
        \end{enumerate}
  \item Dans cette deuxième question, on suppose connues les conditions aux
        limites $V_K = 0$ et $V_A = V_0$, la distance $KA = l$. Le problème
        peut alors se résoudre par la méthode de tir.
        \begin{enumerate}[label=\alph*-]
          \item Expliquer le principe de la méthode de tir.
          \item Écrire alors l'algorithme de résolution de (1) par la méthode
                de tir.
        \end{enumerate}
\end{enumerate}

\subsection*{Exercice 7}

Soit une équation de la forme :
\[
  y''' = f(x, y, y', y'')
  \qquad\text{avec}\quad y(a) = y_1 \quad\text{et}\quad y(b) = y_2 .
\]
On veut résoudre cette équation par la méthode de tir.
\begin{enumerate}[label=\arabic*)-]
  \item Établir le problème aux valeurs initiales résultant.
  \item Le problème aux valeurs initiales résultant doit être résolu par la
        méthode prédicteur-correcteur.
        \begin{enumerate}[label=\alph*)-]
          \item Expliquer le principe de la méthode prédicteur-correcteur sur
                une équation différentielle du premier ordre.
          \item En déduire l'algorithme de résolution du problème
                différentiel du troisième ordre ci-dessus.
        \end{enumerate}
\end{enumerate}

\subsection*{Exercice 8}

\begin{enumerate}[label=\arabic*-]
  \item On considère une fonction $y$ de classe $C^4$ sur un intervalle
        $[a,b]$.
        \begin{enumerate}[label=\alph*-]
          \item En utilisant la formule d'interpolation polynomiale de degré 2,
                établir que $y' = \bigl(y(x+h) - y(x-h)\bigr)/2h$ où $h$ est
                le pas de dérivation.
          \item Déduire de a) une expression pour $y''$.
        \end{enumerate}
  \item Soit le problème aux valeurs aux limites
        $y'' = p(x)y' + q(x)y + g(x)$ défini sur $[a,b]$ avec $y(a) = c$ et
        $y(b) = d$. On veut résoudre ce problème par la méthode des
        différences finies.

        Établir à partir de la question 1) le système algébrique découlant du
        problème.
  \item On veut résoudre le système algébrique ci-dessus par la méthode
        itérative de Jacobi.
        \begin{enumerate}[label=\alph*-]
          \item Expliquer le principe de la méthode itérative de Jacobi.
          \item Écrire alors l'algorithme permettant de résoudre le problème.
          \item Traduire l'algorithme en Pascal.
        \end{enumerate}
\end{enumerate}

\subsection*{Exercice 9}

Une particule de masse $M$ se déplace dans un potentiel de la forme :
\[
  u(y) = \frac{1}{2}ky^2 + \frac{1}{4}k_1y^4 .
\]
Le milieu présente une viscosité linéaire de coefficient $\alpha$ et la
particule est soumise à une force extérieure $f(t) = F\cos\omega t$.
\begin{enumerate}[label=\arabic*)-]
  \item Établir l'équation du mouvement de la particule et adimensionner cette
        équation (rendre l'équation sans dimension).
  \item On donne $y(0) = \dot{y}(0) = 0$ et on veut utiliser la méthode de
        Runge-Kutta d'ordre 4 pour trouver les valeurs de $y$ et $\dot{y}$ en
        fonction du temps $t$.

        Établir l'algorithme de résolution numérique de l'équation
        adimensionnée trouvée à la question (1).
  \item À partir de la résolution numérique, on obtient un ensemble de valeurs
        discrètes $y_i = y(t_i)$, l'expression analytique $y(t)$ étant
        inconnue, on propose la formule approchée suivante :
        \[
          y(t) = \sum_{n=1}^{N}\bigl(a_n\cos n\omega t + b_n\sin n\omega t\bigr).
        \]
        Expliquer comment on peut utiliser la méthode des moindres carrés pour
        déterminer les valeurs des coefficients $a_n$ et $b_n$. Établir les
        algorithmes permettant de déterminer ces coefficients.
\end{enumerate}

\subsection*{Exercice 10}

On considère un système algébrique linéaire $n \times n$ de la forme $AX = b$
où $A = (a_{ij})$, $b = (b_i)$ et $X = (x_i)$ avec $i, j = 1, 2, \dots, n$.

On veut résoudre ce système par la méthode d'élimination de Gauss ou du
pivot.
\begin{enumerate}[label=\arabic*-]
  \item Indiquer le principe de la méthode.
  \item Établir l'algorithme correspondant.
  \item De l'algorithme ci-dessus, on obtient un système triangulaire. Établir
        l'algorithme de résolution du système triangulaire obtenu.
  \item Traduire ce dernier algorithme en Pascal.
\end{enumerate}

\subsection*{Exercice 11}

Le mouvement d'un pendule est décrit par l'équation différentielle
\[
  \frac{\dd^2 Y}{\dd t^2} + \omega_0^2\sin Y = 0
\]
où $\omega_0$ est la pulsation propre du pendule et $Y$ le déplacement
angulaire.
\begin{enumerate}[label=\arabic*-]
  \item On veut résoudre numériquement cette équation par la méthode de
        Runge-Kutta d'ordre 2.
        \begin{enumerate}[label=\alph*-]
          \item Établir l'algorithme de résolution de cette équation.
          \item Traduire cet algorithme en Fortran.
        \end{enumerate}
  \item Lorsque le pendule oscille entre $Y_0$ et $-Y_0$, sa période est
        définie par la relation
        \begin{equation}
          T(Y_0) = \frac{2T_0}{\pi}\int_0^{\pi/2}\frac{\dd x}{\sqrt{1 - A(\sin x)^2}}
          \tag{2}
        \end{equation}
        où $T_0 = 2\pi/\omega_0$ et $A = \sin^2(Y_0/2)$.
        \begin{enumerate}[label=\alph*-]
          \item Après avoir établi la formule composite des trapèzes, donner
                l'algorithme de calcul de l'intégrale (2).
          \item Traduire cet algorithme en Pascal et en Fortran.
        \end{enumerate}
  \item De l'intégration numérique de l'intégrale (2), on obtient le tableau
        suivant :
        \begin{center}
          \begin{tabular}{|l|c|c|c|c|}
            \hline
            $Y_0$ & $\pi/18$ & $\pi/9$ & $\pi/6$ & $2\pi/9$\\ \hline
            $T(Y_0)/T_0$ & 1,002 & 1,008 & 1,020 & 1,030\\ \hline
          \end{tabular}
        \end{center}
        On exprime alors $T(Y_0)$ sous la forme $T(Y_0) = T_0(a + bY_0^2)$. En
        utilisant la méthode des moindres carrés, déterminer les valeurs
        numériques de $a$ et de $b$.
\end{enumerate}

\subsection*{Exercice 12}

On mesure la capacité calorifique $C_v$ d'un corps à volume constant et on
obtient les résultats suivants
\begin{center}
  \begin{tabular}{lccc}
    $T$ (K) & 0.1 & 0.3 & 0.5\\
    $C_v \times 10^4$ cal/K.mole & 0.211 & 0.694 & 1.361
  \end{tabular}
\end{center}
\begin{enumerate}[label=\arabic*-]
  \item On se propose de déterminer la chaleur $Q$ absorbée par une mole du
        corps lorsque sa température passe de 0.1~K à 0.5~K ; calculer $Q$ par
        la formule de Simpson.
  \item En réalité, le tableau ci-dessus est une partie d'un tableau plus
        général contenant de nombreuses valeurs expérimentales. On propose la
        formule $C_v = a + bT$. On veut alors déterminer les coefficients $a$
        et $b$ par la méthode des moindres carrés.
        \begin{enumerate}[label=\alph*-]
          \item Définir le système d'équations vérifié par $a$ et $b$.
          \item Établir des algorithmes permettant d'évaluer numériquement les
                coefficients du système d'équations et de calculer les
                valeurs de $a$ et de $b$.
          \item Traduire ces algorithmes en Pascal.
          \item On veut calculer $Q$ par la formule composite de Simpson à
                partir des valeurs expérimentales de $C_v$. Établir
                l'algorithme de calcul.
        \end{enumerate}
\end{enumerate}

\subsection*{Exercice 13}

On considère l'équation intégrale :
\begin{equation}
  \phi(x) = f(x) + \int_a^b k(x,t)\,\phi(t)\,\dd t
  \tag{1}
\end{equation}
où $k(x,t)$ est une fonction connue. On subdivise le domaine $[a,b]$ en $n$
points $t_i = ih$ équidistants.
\begin{enumerate}[label=\arabic*)-]
  \item L'équation (1) peut être approximée par la formule composite des
        trapèzes. Elle prend alors la forme :
        \begin{equation}
          \phi(x) = f(x) + (b - a)\sum_{i=1}^{n} c_i\,k(x, t_i)\,\phi_i
          \tag{2}
        \end{equation}
        où $\phi_i = \phi(t_i)$.

        Donner les expressions des coefficients $c_i$.
  \item Sachant que l'équation (2) est vérifiée en tout point $x_j = t_j$,
        montrer que l'équation intégrale (1) se réduit au système
        algébrique :
        \begin{equation}
          \phi_j = f_j + (b - a)\sum_{i=1}^{n} c_i\,k(t_j, t_i)\,\phi_i
          \tag{3}
        \end{equation}
        où les $\phi_j$ et $\phi_i$ sont les inconnues.
  \item Le système algébrique (3) peut se mettre sous la forme :
        \begin{equation}
          AX = b
          \tag{4}
        \end{equation}
        où $A$ est une matrice $(n,n)$, $X$ et $b$ sont les matrices colonnes
        de $n$ éléments. Définir ces matrices.
  \item On veut résoudre le système (4) par la méthode de Gauss.
        \begin{enumerate}[label=\alph*)-]
          \item Donner l'algorithme de triangularisation du système.
          \item Écrire l'algorithme de résolution du système triangularisé.
          \item Traduire ce dernier algorithme en Pascal et Fortran.
        \end{enumerate}
  \item $\phi(x)$ possède un zéro dans l'intervalle $[a,b]$ et on se propose
        de trouver ce zéro en utilisant la méthode de bissection. Décrire
        cette méthode et donner son algorithme.
\end{enumerate}

\subsection*{Exercice 14}

À l'instant initial $t = 0$, la température d'un mur d'épaisseur $L$ vaut
$T_0$. Sur la surface $x = 0$, la température varie suivant la loi :
\[
  T(x = 0, t) = T_0 + T_1\sin\frac{\pi}{2}t \qquad\text{pour } t > 0 .
\]
Par un mécanisme approprié, la face $x = L$ est maintenue à la température
constante $T_L$. L'équation de la chaleur dans le milieu est alors :
\[
  \frac{\partial T}{\partial t} = \alpha\frac{\partial^2 T}{\partial x^2}
  \qquad\text{où } \alpha \text{ est la diffusibilité.}
\]
\begin{enumerate}[label=\arabic*)-]
  \item On se propose de résoudre numériquement le problème ainsi posé par la
        méthode des différences finies progressives.
        \begin{enumerate}[label=\alph*)-]
          \item Donner la forme discrète du problème. Quel est l'ordre de
                l'erreur de discrétisation ?
          \item Écrire l'algorithme de résolution du système discret obtenu.
          \item Traduire cet algorithme en Pascal.
        \end{enumerate}
  \item On veut maintenant utiliser la méthode de Richardson pour résoudre le
        même problème.
        \begin{enumerate}[label=\alph*)-]
          \item Donner le nouveau système discret et indiquer l'ordre et
                l'erreur de discrétisation.
          \item En admettant que l'on utilise le schéma discret de la question
                (1) pour approximer les températures
                $T_{i,1} = T(x = i\Delta_x, \Delta t)$, établir l'algorithme
                permettant de trouver la température aux instants ultérieurs.
          \item Traduire cet algorithme en Fortran.
        \end{enumerate}
\end{enumerate}

\subsection*{Exercice 15}

\begin{enumerate}[label=\arabic*)]
  \item On considère l'équation de la chaleur
        $\dfrac{\partial u}{\partial t} = a^2\dfrac{\partial^2 u}{\partial x^2}$
        définie sur $D = \R_+ \times\, ]0, l[$ avec
        $u(0,t) = u(l,t) = 0$, $u(x,0) = f(x)$.

        On veut résoudre cette équation par le schéma explicite avec
        $\dfrac{\partial u}{\partial t} = \dfrac{u_{i,j+1} - u_{i,j}}{k}$.

        Établir que la condition de stabilité de ce schéma est
        $\dfrac{a^2k}{h^2} \le \dfrac{1}{2}$ où $k$ et $h$ sont les pas
        temporel et spatial.
  \item Que devient cette condition de stabilité dans le cas d'un problème
        plan : $\dfrac{\partial u}{\partial t} = a^2\Delta u$ ? On prendra
        $h_1$ et $h_2$ les pas spatiaux.
  \item On considère l'équation d'advection-diffusion dans le plan :
        \[
          \frac{\partial\Omega}{\partial t} + u_1\frac{\partial\Omega}{\partial x}
          + u_2\frac{\partial\Omega}{\partial y} = \nu\Delta\Omega .
        \]
        On la discrétise par le schéma centré en espace. Établir que la
        condition de stabilité est $\dfrac{\nu k}{h^2} \le \dfrac{1}{4}$ où
        $k$ est le pas temporel et $h$ est le pas spatial suivant $x$ et sur
        $y$.

        Que devient cette condition dans le cas unidimensionnel ?
\end{enumerate}

\subsection*{Exercice 16}

Soit l'équation d'onde suivante :
\[
  \frac{\partial^2 u}{\partial t^2} + a\frac{\partial u}{\partial t}
  = c^2\frac{\partial^2 u}{\partial x^2} + du^3
\]
où $(t, x) \in \R_+ \times\, ]0, l[$ et $u(0,t) = u(l,t) = 0$,
$u(x,0) = f(x)$, $\dfrac{\partial u}{\partial t}(x,0) = g(x)$.

$a$, $c$ et $d$ sont des constantes.
\begin{enumerate}[label=\arabic*)-]
  \item Définir chaque terme de cette équation.
  \item On veut résoudre numériquement cette équation par les méthodes de
        différences finies centrées. Donner la forme discrète du problème
        ainsi posé.
  \item Exprimer $u_{i,2} = u(x = ih, t = 2k)$ en fonction des valeurs des
        fonctions $f$ et $g$ aux points $i-2$ ; $i-1$ ; $i$ ; $i+1$ et $i+2$
        et des pas $h$ et $k$.
  \item Établir un algorithme de résolution du système discret obtenu à la
        question (2).
  \item Traduire l'algorithme dans un langage de programmation de votre choix.
\end{enumerate}

\subsection*{Exercice 17}

\begin{center}
  \begin{tikzpicture}[>=Stealth]
    % coupe du tube (1 cm pour 125 mm) : région hachurée entre les deux rectangles
    \fill[pattern=north east lines, even odd rule]
      (0,0) rectangle (8,6.4) (2,1.6) rectangle (6,4.8);
    \draw[thick] (0,0) rectangle (8,6.4);
    \draw[thick] (2,1.6) rectangle (6,4.8);
    \node at (3.2,3.6) {$p_s$};
    % cotes intérieures
    \draw[<->, dashed] (2,2.4) -- (6,2.4) node[pos=0.5, below] {$L_1$};
    \draw[<->, dashed] (5.2,1.6) -- (5.2,4.8) node[pos=0.6, left] {$l_1$};
    % cotes extérieures
    \draw[<->] (-0.5,0) -- (-0.5,6.4) node[midway, left] {$l$};
    \draw[<->] (0,-0.5) -- (8,-0.5) node[midway, below] {$L$};
  \end{tikzpicture}
\end{center}
\[
  L = 1000~\text{mm}, \quad L_1 = 500~\text{mm}, \quad l = 800~\text{mm}, \quad l_1 = 400~\text{mm}
\]
Le dispositif représente la coupe d'un tube rectangulaire dans lequel
s'écoule dans sa partie centrale un liquide avec une pression $p_s$. La
pression sur la surface extérieure est nulle. À chaque point de la région
hachurée, la distribution des pressions vérifie l'équation
\[
  \frac{\partial^2 p}{\partial x^2} + \frac{\partial^2 p}{\partial y^2} = 0 .
\]
\begin{enumerate}[label=\arabic*)-]
  \item Faites un adimensionnement de cette équation et précisez les nouvelles
        conditions sur la frontière. La pression adimensionnée est $u$.
  \item \begin{enumerate}[label=\alph*)-]
          \item En utilisant la méthode des différences finies, établir
                l'équation discrète vérifiée par $u_{i,j}$ en tout point
                $(x_i = ih, y_j = jk)$.
          \item Discrétiser les conditions aux limites.
        \end{enumerate}
  \item On se propose de trouver les $u_{i,j}$ par la méthode itérative de
        Gauss-Seidel.
        \begin{enumerate}[label=\alph*)-]
          \item Décrire cette méthode itérative.
          \item Établir l'algorithme de calcul des $u_{i,j}$.
        \end{enumerate}
  \item On suppose que le tube a maintenant une forme cylindrique de rayon
        intérieur $R_1 = 400$~mm et de rayon extérieur $R_2 = 800$~mm.
        \begin{enumerate}[label=\alph*)-]
          \item Quelle est alors l'équation vérifiée par la pression $p$ ou la
                grandeur adimensionnée $u$ ?
          \item Proposer une méthode de discrétisation de cette nouvelle
                équation.
        \end{enumerate}
\end{enumerate}

\subsection*{Exercice 18}

\begin{center}
  \begin{tikzpicture}[scale=0.6,>=Stealth]
    % plaque L x l avec une encoche centrale de largeur L/3 et de hauteur l/2
    \fill[pattern=north east lines]
      (0,0) -- (3,0) -- (3,3) -- (6,3) -- (6,0) -- (9,0) -- (9,6) -- (0,6) -- cycle;
    \draw[thick]
      (0,0) -- (3,0) -- (3,3) -- (6,3) -- (6,0) -- (9,0) -- (9,6) -- (0,6) -- cycle;
    % cotes
    \draw[<->] (0,6.6) -- (9,6.6) node[midway, above] {$L$};
    \draw[<->] (-0.6,0) -- (-0.6,6) node[midway, left] {$l$};
    \draw[<->] (0,-0.6) -- (3,-0.6) node[midway, below] {$\frac{L}{3}$};
    \draw[<->] (6,-0.6) -- (9,-0.6) node[midway, below] {$\frac{L}{3}$};
    \draw[<->] (5.5,0) -- (5.5,3) node[midway, left] {$\frac{l}{2}$};
  \end{tikzpicture}
\end{center}

La répartition d'une grandeur physique $u$ sur la plaque représentée vérifie
l'équation de Poisson suivante
\[
  \frac{\partial^2 u}{\partial x^2} + \frac{\partial^2 u}{\partial y^2} = f(x, y)
\]
avec $u = g(x, y)$ sur la frontière.
\begin{enumerate}[label=\arabic*)-]
  \item En utilisant la méthode des différences finies, établir l'équation
        discrète vérifiée par $u_{i,j}$.
  \item On se propose de trouver les $u_{i,j}$ par la méthode itérative de
        Jacobi.
        \begin{enumerate}[label=\alph*)-]
          \item Décrire cette méthode.
          \item Établir l'algorithme de calcul de $u_{i,j}$.
        \end{enumerate}
\end{enumerate}

\subsection*{Exercice 19}

\begin{center}
  \begin{tikzpicture}[scale=0.8,>=Stealth]
    \fill[pattern=north east lines] (0,0) -- (4,0) -- (0,4) -- cycle;
    \draw[thick] (0,0) -- (4,0) -- (0,4) -- cycle;
    \draw[<->] (-0.5,0) -- (-0.5,4) node[midway, left] {$h$};
    \draw[<->] (0,-0.5) -- (4,-0.5) node[midway, below] {$l$};
  \end{tikzpicture}
\end{center}

La répartition d'une grandeur physique sur la plaque triangulaire ci-dessus
vérifie l'équation de Laplace :
$\dfrac{\partial^2 u}{\partial x^2} + \dfrac{\partial^2 u}{\partial y^2} = 0$
avec $u = 0$ sur l'hypoténuse et $\dfrac{\partial u}{\partial n} = a$ sur les
autres côtés ($n$ est la normale par rapport aux côtés considérés).
\begin{enumerate}[label=\arabic*)-]
  \item Discrétiser ce problème par la méthode des différences finies.
  \item Établir l'algorithme de calcul des valeurs discrètes $u_{i,j}$ sur
        toute la plaque.
\end{enumerate}
